\section{Limitations}\label{sec:limits}

\subsection{Threats to validity}

\textbf{Internal validity.} The seeds are few: \nSeedsOne{} for the crosstalk device, \nSeedsOneFull{} for the full device and \schroffXN{} for each loss, so the intervals are wide and no significance test is reported; the equivalence of the two losses is an agreement within seed ranges and not a test of equivalence. The benchmark cells at \benchOneBudget{} iterations hold \benchOneDevSeeds{} seeds for the refinement, so the geometry results at that budget are ranges, and the equal-time rows, which hold \benchOneWinMinSeeds{} seeds on each side, are read with a rule that compares the refinement with phase-fixed WGS only. The normalized-error loss is our implementation of the loss of Schroff et al.\ on point traps with a closed-form gradient: it is not their code, their automatic differentiation, their continuous targets convolved with the point-spread function or their camera feedback, and a result of the paper about that loss is a result about this implementation. The variants of the ablation are re-implementations of the loop of the module and were verified to reproduce it, so the ablated rows are not a different method. The times are wall-clock seconds of one machine on which other jobs also ran during parts of the measurements, so they are read as orders of magnitude and not as ratios of a quiet machine.

\textbf{External validity.} Every result is simulated. The crosstalk parameters come from a published fit on another device; the source gives the width in inverse pixels and does not define the unit of the frequency further, and the engine's reading, cycles per SLM pixel on the upsampled grid, is an assumption of the implementation that was not checked against the source's code. The \parBits{}-bit quantizer, the fill factor and the flicker have no published fit in this model, and random flicker is scored and not modelled by the optimizer. An instrument with another kernel would give another factor. The tasks are lattices of equal traps on bins at one wavelength and one pitch, a Gaussian beam and no aberrations; the benchmark adds other geometries on one device, and the three-dimensional and shaped-field targets are outside the scope of the method.

\textbf{Construct validity.} The spread and the efficiency are model quantities, and the factor measures them against phase-fixed WGS without camera feedback, whereas the experiments close the gap with feedback; a method that is two orders of magnitude flatter before feedback may or may not need fewer feedback steps, and that was not measured. The factor on the lattice task is large partly because the baseline's spread, set by the crosstalk, does not fall with more iterations (the baseline at 100 iterations is no better than at 30), so the size of the factor is a property of this device model and of this lattice and not a general figure.

\subsection{Scope and what is not claimed}

The paper does not claim a new loss, a new device model or the idea of a crosstalk model inside a conjugate-gradient minimization of a trap array: these belong to Schroff et al., and the loss they use is run here as the baseline and matches the default. Neither the refinement nor the loss is shown to improve a camera-corrected experiment or the survival of atoms. The method is built for traps on FFT bins at the focal plane without vortices and costs about $u^2$ times a plain iteration per step and, in the measurements, minutes to tens of minutes per hologram on a CPU, so it is baked and not live. The work of Persson et al.\ is known only through its abstract, so the paper does not claim that no earlier solver combines a crosstalk model with the weighted-GS weights. The paper also does not test the sensitivity of the refinement to a kernel that differs from the one the device displays, which is the case in any real instrument whose kernel is fitted with error; a hologram optimized on a wrong kernel would carry part of the crosstalk error that the method removes in the model. The moving-trap sequences of the studio are not part of this paper.

\section{Reproducibility}\label{sec:repro}

The code, the data and the committed artifacts are in the repository \url{https://github.com/fsantibanezleal/CEFOP_DinHot} at release \engineRelease{} (Apache-2.0); the studio runs at \url{https://dinhot.fasl-work.com}. Every number, table and figure of the paper is produced by three scripts from the repository root. The measurement script \path{manuscripts/device-aware-flicker-bounded/data/measure.py} runs the methods of the engine and writes
\texttt{beyond.json}, the macro file read by the text and the result tables; with no argument it measures every part (the main measurements took about \mainMinutes{} minutes and the ablation about \ablSeconds{} minutes on the development machine, CPU only, with five seeds on each device). The comparison of the weight update and of the normalized-error loss on the benchmark lattices (Table~\ref{tab:ablb1geo}) is written by \path{manuscripts/device-aware-flicker-bounded/data/b1_weights_geo.py}, one process per geometry (\texttt{run G1} to \texttt{run G4}, about half an hour each or more) and then \texttt{write}. The figure script \path{manuscripts/device-aware-flicker-bounded/figures/make_figs.py} writes every figure as a
vector PDF from \texttt{beyond.json} and the geometry files, and writes \texttt{fig\_numbers.tex}, the macros of the derived numbers of Section~\ref{sec:results} (times, budgets, counts) and the check of Eq.~\ref{eq:losses}, which compares the engine's two cost classes on a small array behind the crosstalk device: the two values agree to \lossValueDiff{} and the gradient differs from the prediction of the proposition by \lossGradDiff{} in the largest component. The document builds with three passes of \texttt{pdflatex} from \texttt{tex/}. The benchmark file and the committed manifests that the scripts read are in \texttt{data/derived/}. The acceptance tests of the method are \texttt{test\_b1\_device\_wgs.py}, with the tests of the loss option and of the default, and the ablated variants are checked against the module inside the measurement script. All methods are deterministic given a seed, so a rerun reproduces the numbers up to the floating point of the machine's FFT; the times differ between runs.

\section{Conclusions}\label{sec:conclusions}

The device-aware conjugate-gradient refinement does not improve on the normalized-error loss of Schroff et al.\ through the same device model and conjugate gradient, and the paper reports that as a negative result. On a $10\times10$ lattice behind a crosstalk device the normalized-error loss reaches \schroffXSig{} percent against \ablBFullSig{} percent for the efficiency-level loss, and on the full device \schroffFSig{} against \ablBFFullSig{} percent, where it has the lower median, within the ranges of the seeds on both. The safeguard did not act on those seeds, and the multiplicative weight update of weighted GS makes the spread \ablBWeightsRatio{} times higher on the acceptance task and \ablBFWeightsRatio{} times higher on the full device, so it is off in the method. What lowers the spread is the device model in the loop, which Schroff et al.\ already use: it brings a geometric-mean factor of \ratioOneGm{} (95 percent interval \ratioOneCiLo{} to \ratioOneCiHi{}) over phase-fixed weighted GS at an efficiency ratio of at least \effRatioOneMin{}, and the gain disappears when the model is removed from the loop. The gain does not carry to every geometry and it is not free: on the benchmark's device with flicker the method meets both thresholds on the ring, the random set and the two large sets at 100 iterations, with two seeds per cell, and falls short on the square lattice and the quasi-periodic set; at an equal wall-clock time it is read as better than phase-fixed WGS on \benchOneDevTimeProposed{} of \benchOneDevCells{} geometries and ties on \benchOneDevTimeTie{}; and it takes \benchTenSecRefMin{} to \benchTenSecRefMax{} s per hologram at 10 outer iterations against \benchTenSecMsevenMin{} to \benchTenSecMsevenMax{} s for the baseline, and \benchSecGeoOneBOne{} s against \benchSecGeoOneMseven{} s on the square lattice at the 100-iteration row.
The next steps are the measurement on a device with a fitted model and a camera in the loop, the sensitivity of the refinement to a kernel that differs from the displayed one, and the reading of the work of Persson et al.\ in full to place the method among the solvers that already carry a crosstalk model.
